\documentclass[11pt]{article}
\usepackage[margin=1.05in]{geometry}
\usepackage{amsmath,amsthm,amssymb,mathtools,bm}
\usepackage{booktabs,array,enumitem}\usepackage[expansion=false]{microtype}
\usepackage{xcolor}
\usepackage[colorlinks=true,linkcolor=blue!60!black,citecolor=blue!60!black,urlcolor=blue!60!black]{hyperref}
\theoremstyle{plain}
\newtheorem{theorem}{Theorem}[section]\newtheorem{proposition}[theorem]{Proposition}\newtheorem{lemma}[theorem]{Lemma}
\newtheorem{corollary}[theorem]{Corollary}\newtheorem{conjecture}[theorem]{Conjecture}
\theoremstyle{definition}\newtheorem{definition}[theorem]{Definition}\newtheorem{remark}[theorem]{Remark}
\newtheorem{example}[theorem]{Example}
\newcommand{\R}{\mathbb R}\newcommand{\Z}{\mathbb Z}\newcommand{\E}{\mathbb E}\newcommand{\Var}{\operatorname{Var}}\newcommand{\Cov}{\operatorname{Cov}}
\newcommand{\relu}{\operatorname{relu}}\newcommand{\diag}{\operatorname{diag}}\newcommand{\tr}{\operatorname{tr}}
\newcommand{\Tr}{\operatorname{Tr}}\newcommand{\1}{\mathbf 1}\newcommand{\cA}{\mathcal A}\newcommand{\cB}{\mathcal B}
\newcommand{\cF}{\mathcal F}\newcommand{\cT}{\mathcal T}\newcommand{\cH}{\mathcal H}\newcommand{\aff}{\mathfrak{aff}}
\newcommand{\He}{\mathrm{He}}\newcommand{\Fa}{\Phi_{\mathrm F}}\newcommand{\cZ}{\mathcal Z}\newcommand{\Sym}{\operatorname{Sym}}
\newcommand{\Aff}{\mathrm{Aff}}\newcommand{\Ryshkov}{\mathcal R}
\newcommand{\Ap}{\mathsf A}\newcommand{\J}{\mathsf J}\newcommand{\cG}{\mathcal G}\newcommand{\Br}{\mathsf{Br}}\newcommand{\Pa}{\mathsf P}
\title{\textsc{Localization in the commutant, stage 16:}\\ \large the commutant tower, operator growth, and transport memory.\\
What the companion notes establish, what the network forces,\\ and where an estimator is actually free}
\author{}
\date{10 October 2026}
\begin{document}
\maketitle
\begin{abstract}
This stage starts from two companion notes and from the experimental session's tests of stage~13, all supplied by the user. The notes are
\emph{Branching, angular geometry and joint dynamics} and \emph{Wick geometry, covariance holonomy and the closure of Kikuchi dynamics}.
\begin{itemize}[leftmargin=*]
\item We audit the notes' claims numerically, all of which hold.
\item We prove what in them is forced for the network, and we measure what decides between design directions.
\item We test two non-standard estimator designs that the notes suggest; both fail, and we say exactly why.
\end{itemize}
\begin{itemize}[leftmargin=*]
\item \emph{The commutant tower (proved; checked).} The weight ensemble's symmetry $O(n)$ has the Brauer algebra (pairings) as its
commutant on tensors. The ReLU fold's symmetry $B_n$ has even-block partitions, and hidden-unit relabelling $S_n$ has all partitions. For
$4$ and $6$ legs the dimensions are $3<4<15$ and $15<31<203$.
\begin{itemize}
\item Annealed pre-activation statistics are pairings: Gaussian plus dilation, exactly the companion's mixture formula
$\kappa_4=\Cov(Q)$. In particular annealed odd cumulants of pre-activations vanish.
\item Folds create the other blocks. The ``bulk three-site structure'' left in v56 is the quenched three-block sector.
\item The hard-core normalizer theorem says that the hyperoctahedral group, $B_n=\mathrm{Aut}(\Z^n)\cap O(n)$, is the only
transport keeping neuron-diagonal data closed. Gaussian weights are almost surely not in it, so the slices must be regenerated from transported
legs: \emph{transport memory} is forced.
\end{itemize}
\item \emph{Operator growth (measured on official network~0).} As functions of the input, the neurons' first-chaos share falls from
$73\%$ (layer~1) to $19\%$ (layer~16). The mean degree of their higher chaos grows linearly, $2.8\to17.1$, about one degree per layer. A
fixed-basis (Pauli-path, input-Hermite) expansion is therefore hopeless at depth. The layer-wise Wick re-centring of
Ebrahimi-Fard--Patras--Tapia--Zambotti, a frame that co-evolves with the instance's weights, is what makes the problem low-order. The
interdependence of representation and weights is not optional; it is forced.
\item \emph{Where the correction lives (measured).}
\begin{itemize}
\item $60$--$66\%$ of what any estimator must add to the Gaussian closure is an isotropic function of each row's field (the Brauer and
dilation sector). The rest is anisotropic.
\item The exact first-order chain captures $99.1$--$99.7\%$ of the correction.
\item Its residual is incoherent and born late, consistent with the experimental session's finding that $0.4$--$3.7\%$ of v56's error lies
along the mean.
\end{itemize}
\item \emph{Two designs tested and ruled out.}
\begin{itemize}
\item The leaf (wall-sum) estimator, using the companion's Dirac energies (carré du champ), fails by a factor $10^5$ already at layer~2.
The global Dirac energy exceeds the transverse one by the Poincaré deficit, the higher-chaos energy, and the exact identity relies on
cancellations between layers.
\item Minimal-sensitivity Wick frames do not beat the centred (BPHZ) frame.
\end{itemize}
\item \emph{What is free.} We state a row-exchangeable self-calibration theorem: an estimator can learn its own correction from its own
weights. We then give the design space that remains: frame family, diagram blocks, rows, layers, instruments.
\end{itemize}
Checks: \texttt{scripts/verify\_stage16.py}.
\end{abstract}
\tableofcontents
\section{The picture, and an accounting of stage 15}
\label{sec:summary}
\subsection{An accounting of stage 15}
Stage~15 refused to assume that an estimator is a law propagator, and derived the error cocycle from first principles. Its proposed
``theory-native estimator'' was still chain-shaped, however. It kept an exact backbone, transported births and allocated effort; the
Heisenberg reads and the knapsack are reorganisations of the same object. This stage asks the prior question: which features of a chain are
\emph{forced} by the mathematics, and which are free? Every claim below is labelled forced (a theorem or a measurement that leaves no choice)
or free (a design freedom), and the free ones are tested where possible.

\subsection{What is forced}
\begin{enumerate}[leftmargin=*]
\item \textbf{Co-evolving frames} (Section~\ref{sec:growth}; measured). Network~0's neurons, as functions of the input, have mean
higher-chaos degree $2.8$ at layer~1 and $17.1$ at layer~16. A representation in a fixed basis is hopeless. The representation must be
re-centred at every layer by the instance's own law, the Wick map $\mu_l^{-1}\star\mathrm{id}$. The estimator's frame therefore depends on
the weights dynamically.
\item \textbf{Algebraic sectors} (Section~\ref{sec:commutant}; proved). Annealed pre-activation statistics are Brauer elements (pairings),
which covers the Gaussian and the dilation. The folds create every other partition block. The instance-specific (quenched) parts are
orthogonal to all annealed ones.
\item \textbf{Transport memory} (Section~\ref{sec:transport}; proved). No closure on neuron-diagonal data alone can be exact under a
transport outside the hyperoctahedral group $B_n=\mathrm{Aut}(\Z^n)\cap O(n)$. Gaussian weights are almost surely outside it. Births must be
carried with their own transported legs, or recomputed.
\item \textbf{Incoherence of the residual} (stage~15; confirmed by the experimental session). Once the modular sector is right, the
error is a sum of independent row-births, and no global correction reaches it.
\end{enumerate}

\subsection{What is free, and what testing it showed}
\begin{center}\small
\begin{tabular}{>{\raggedright\arraybackslash}p{4.6cm}>{\raggedright\arraybackslash}p{5.0cm}>{\raggedright\arraybackslash}p{4.4cm}}\toprule
freedom & candidate & status\\\midrule
frame parameters & minimal-sensitivity (PMS) Gaussian frames & no gain over centring (checked)\\
frame family & non-Gaussian (mixture, skew) frames & open\\
decomposition of the mean & leaf/wall sum with Dirac energies & fails without transverse conditioning (checked)\\
which blocks to evaluate & partition-algebra diagram selection by annealed variance & framework; open\\
which rows & row-exchangeable self-calibration & theorem; protocol for test\\
which layers & cocycle allocation (late layers) & measured value; untested in a system\\
instruments & task-preserving branching on collective walls & theory; cost doubles per branch\\\bottomrule
\end{tabular}
\end{center}

\subsection{The numbers that decide}
\begin{center}\small
\begin{tabular}{>{\raggedright\arraybackslash}p{8.4cm}>{\raggedright\arraybackslash}p{5.2cm}}\toprule
quantity (official network~0) & value\\\midrule
first-chaos share of neurons, layer $1\to16$ & $0.73\to0.19$\\
mean degree of higher chaos, layer $1\to16$ & $2.8\to17.1$ (about $+1$ per layer)\\
isotropic (field) share of truth $-$ Gaussian closure, layers $4$--$16$ & $0.60$--$0.66$\\
share of truth $-$ Gaussian closure captured by the exact first-order chain & $0.991$--$0.997$\\
wall-sum estimator with global Dirac energies, layer 2 & MSE $3.5\times10^{-2}$ (Gaussian closure $1.5\times10^{-7}$)\\
Dirac energy / variance of pre-activations, layers $1,2,4,16$ & $1.00$, $1.47$, $2.54$, $13.7$\\\bottomrule
\end{tabular}
\end{center}

\section{Audit of the companion notes}
\label{sec:audit}
\subsection{Branching, angular geometry and joint dynamics}
The note reads Zhou's paper constructively. We checked its finite statements and kept the following.
\begin{itemize}[leftmargin=*]
\item \emph{Poisson tensor growth.} $\kappa_t=e^{t(\chi_V-d)}$ is a Poisson mixture of normalized characters of $V^{\otimes k}$. The label chain
$K(\lambda,\nu)=c^\nu_{\lambda V}\dim V_\nu/(d\dim V_\lambda)$ is stochastic. For spin $\frac12$ appended to spin $j$ the branches are
$p_\pm=(j+1)/(2j+1),\,j/(2j+1)$ (checked at $j=\frac12,1,\frac72$).
\item \emph{Multiplicity as a logical qubit.} In the three-spin decoherence-free subsystem with the stated basis, $S_{12}=-Z$,
$S_{23}=\frac12Z+\frac{\sqrt3}2X$ and $[S_{12},S_{23}]=-i\sqrt3\,Y$ (checked exactly). Appending a spin preserves the multiplicity register even
conditionally on the branch, because each branch effect is scalar on $V_j$ (Schur).
\item \emph{The angular tower.} $|\psi_\theta\rangle=(|0\rangle+\theta)/\sqrt2$ gives $\Tr(R_{W,m}A_{W,m})=\int f_W\,d\sigma$ and $T_m^*B_{m-1}=B_m$. These are exact,
including for coupled registers. The isotropic reference is $G_1=\operatorname{diag}(\frac12,\frac1{2n},\dots)$ (checked: $0.5$, $0.0998$--$0.1004$
at $n=5$), faithful and non-tracial, with modular phases $n^{iu}$ on $|0\rangle\langle i|$.
\item \emph{The correction to Zhou v1, Theorem~4.3.} For the signed Poisson walk, $\E[Z_a^2Z_b^2]=ab+2a^2+a$. Monte Carlo gives $3.021$ against
$3.010$; the product of separate means would be $1.33$. The multi-time formula must be read as nested conditional expectations. Stage~15
cited that theorem only for the restriction mechanism, which stands.
\item \emph{The product-defect identity} $P_t(AB)-P_tAP_tB=\int_0^tP_s\Gamma(P_{t-s}A,P_{t-s}B)\,ds$, the Dirac realization
$\|[D_G,\pi(A)]\Omega_\rho\|^2=\Tr(GC_\rho(A))$, moving contacts, and protected logical observables. We use the first two in
Section~\ref{sec:gamma}.
\end{itemize}
What does not transfer by itself:
\begin{itemize}[leftmargin=*]
\item The angular representation needs the cell integrals, so it is exact but not cheap.
\item Growth needs a real increment.
\item The note says both of these itself.
\end{itemize}

\subsection{Wick geometry, covariance holonomy and Kikuchi closure}
\begin{itemize}[leftmargin=*]
\item \emph{Holonomy.} For the symmetric-squeeze connection $\dot T=\frac12\dot QQ^{-1}T$ around a contact-preserving square of side $a$ in
$Q=\begin{psmallmatrix}1&v\\v&1+u\end{psmallmatrix}$, the transport is exactly orthogonal ($|TT^\top-I|\le5\times10^{-13}$). Its angle over $a^2/4$ is
$0.985$, $0.965$, $0.935$ at $a=0.02$, $0.05$, $0.1$, so the leading coefficient $\frac14$ is confirmed with an $O(a^3)$ correction.
\item \emph{Hard-core normalizer.} The coefficient of $z_i^2$ in $(Rz)_j(Rz)_k$ is exactly $0$ for a signed permutation and $0.39$ for a
random rotation.
\item \emph{Mixture fourth cumulant.} $\kappa_4(X)=\Cov(Q_{ij},Q_{kl})+\Cov(Q_{ik},Q_{jl})+\Cov(Q_{il},Q_{jk})$. Monte Carlo gives $1.573$, $0.538$,
$-0.007$ against $1.558$, $0.518$, $-0.001$.
\item \emph{Realization relations.} For $B\sim\mathrm{Bern}(p)$, $b^2-b=0$ but $:\!b^2\!:-:\!b\!:\;=-2p(B-p)$ (checked). Our gates satisfy $g^2=g$, so
a Wick deformation must respect this relation. That is one reason the centred Gaussian frame works on pre-activations, where it is
nowhere violated, and not on gates.
\item \emph{The all-degree sparsifier and the Poisson-coupled weights.} They concern a mean-estimation \emph{dynamics} (sampling), so they
are bounded by stage~9's sampling law; we do not use them as estimators. The convex variational principle with equalized edge energies is
the exact form of stage~15's knapsack (Section~\ref{sec:interdep}).
\end{itemize}

\section{The commutant tower: Brauer, even partitions, partitions}
\label{sec:commutant}
\subsection{Three symmetry groups and their tensor commutants}
Three groups act on a hidden layer's $\R^n$.
\begin{itemize}[leftmargin=*]
\item $O(n)$, under which the Gaussian weight ensemble is invariant on each side.
\item The hyperoctahedral group $B_n$ of signed permutations, which commutes with the fold $|\cdot|$.
\item $S_n$, the relabelling of hidden units, which commutes with $\relu$.
\end{itemize}
\begin{theorem}[the commutant tower; classical, checked]\label{thm:tower}
For $n\ge2k$, the algebras of operators on $(\R^n)^{\otimes k}$ commuting with the diagonal action are as follows.
\begin{itemize}[leftmargin=*]
\item $O(n)$: the Brauer algebra $\Br_k(n)$, spanned by perfect matchings of the $2k$ legs.
\item $B_n$: the span of set partitions of the $2k$ legs with all blocks of even size.
\item $S_n$: the partition algebra $\Pa_k(n)$, spanned by all set partitions.
\end{itemize}
So $\Br_k\subset\Pa^{\rm even}_k\subset\Pa_k$.
\end{theorem}
\begin{proof}
An operator commuting with all coordinate permutations has entries that depend only on the equality pattern of its $2k$ indices, a set
partition. Commuting also with the sign flips forces every block to contain an even number of legs. The $O(n)$ case is Brauer's
theorem.
\end{proof}
\emph{Check.} Invariant dimensions are $\frac1{|G|}\sum_g\chi(g)^{2k}$ for the finite groups and $\E\,(\Tr O)^{2k}$ for Haar $O(n)$.
\begin{itemize}[leftmargin=*]
\item $2k=4$: $S_4$ gives $15$, $B_4$ gives $4$, and $O(6)$ gives $3.01\pm0.02$. The counts are $15$ (Bell), $4$ (even
partitions) and $3$ (pairings).
\item $2k=6$: $203$, $31$, and $15.17\pm0.22$; the counts are $203$, $31$, $15$.
\end{itemize}

\subsection{Which statistics live where}
\begin{proposition}[annealed statistics; proved]\label{prop:annealed}
Let $W_l$ have i.i.d.\ $N(0,\sigma^2)$ entries independent of the upstream state, and $z_l=W_lh_{l-1}$.
\begin{enumerate}[leftmargin=*]
\item $\E_{W_l}[\kappa_k(z_l)]=\sigma^k\sum_\pi\prod_{(a,b)\in\pi}\delta_{i_ai_b}\cdot\langle\kappa_k(h_{l-1}),\pi\rangle$, summed over perfect matchings $\pi$ of the $k$ legs. Annealed
pre-activation cumulants are Brauer elements; in particular \emph{all annealed odd cumulants of pre-activations vanish}.
\item Post-activation annealed statistics are $S_n$-invariant, hence partition-algebra elements. Odd blocks appear because $\relu$
breaks the sign symmetry. Even non-pairing blocks such as $\delta_{iiii}$ come from the fold and odd blocks such as $\delta_{iii}$ from the
asymmetry.
\item If $z\mid Q\sim N(0,Q)$, then $\kappa_4=\Cov(Q_{ij},Q_{kl})+{}$two permutations, a pure pairing element. With $Q=qI$ (the dilation) it is
$\Var(q)(\delta\delta+\delta\delta+\delta\delta)$. Upper-geometry fluctuations create only Brauer-type cumulants. The dilation package of stage~13 is the
Brauer $\kappa_4$.
\end{enumerate}
\end{proposition}
\begin{proof}
(1) Gaussian Wick calculus for $E[W_{i_1j_1}\cdots W_{i_kj_k}]$ pairs the factors, identifying row and column indices together. (2) Relabelling
of hidden units is a symmetry of the ensemble of networks, while sign flips are not. (3) Condition on $Q$ and subtract the three pair
products (checked in Section~\ref{sec:audit}).
\end{proof}
\begin{corollary}[where the residual lives]
Stage~15's isotropic sector is the Brauer span. Three-block content of pre-activations, the ``stars'', has zero annealed part, so it is
\emph{purely quenched}: no counterterm, mean-field tier or annealed correction can reach it. Only contractions with the instance's
transported legs can. This is the experimental session's ``bulk three-site structure''. The dilation, by contrast, is Brauer, and the
fitted counterterms had already absorbed it; that is why its exact package did not pay on v56.
\end{corollary}

\subsection{The noiseless subsystem of the ensemble}
By Schur--Weyl--Brauer duality, $(\R^n)^{\otimes k}=\bigoplus_\lambda V^{O(n)}_\lambda\otimes M_\lambda$, where $\Br_k$ acts on the multiplicity spaces $M_\lambda$. The
ensemble average over $O(n)$ is the twirl that keeps only the multiplicity (leg-pattern, replica-overlap) information. This is the
companion's ``information behind the label'' with the roles reversed. The collective noise here is the randomness of the weights; it
\emph{decoheres the spatial factor}, where the instance-specific information lives, and preserves the multiplicity factor, which is the
annealed (NNGP-type) theory. Annealed theory is the dynamics of the noiseless subsystem. Estimation of one instance is the recovery of
what that twirl destroys, and a quantum Schur transform would not help with it: it gives access to multiplicities, which are classically
cheap here.

\section{Operator growth and the necessity of co-evolving frames}
\label{sec:growth}
\subsection{The chaos ladder of a real network}
Regard a neuron $h(x)$ as a function of the Gaussian input and expand it in Wiener chaos, $h=\sum_kh_k$. Three energies are accessible by
forward-mode differentiation alone:
\[
F=|\E\nabla h|^2=\|h_1\|^2,\qquad V=\Var h=\sum_{k\ge1}\|h_k\|^2,\qquad G=\E|\nabla h|^2=\sum_kk\|h_k\|^2 .
\]
From them we get the higher-chaos share $1-F/V$ and the mean degree of the higher chaos, $\langle k\rangle_{\ge2}=1+(G-V)/(V-F)$. Hutchinson's
identities give unbiased estimates,
\[G=\E_{x,v}(\nabla h\cdot v)^2,\qquad F=\E_{x,x',v}(\nabla h(x)\cdot v)(\nabla h(x')\cdot v),\]
with independent $x,x'$ and a shared probe $v$. On official network~0, with $49152$ inputs and totals over neurons:
\begin{center}\footnotesize\setlength{\tabcolsep}{3pt}
\begin{tabular}{lcccccccccccccccc}\toprule
layer&1&2&3&4&5&6&7&8&9&10&11&12&13&14&15&16\\\midrule
$F/V$&.73&.60&.51&.44&.39&.36&.33&.29&.27&.26&.24&.23&.22&.21&.20&.19\\
$\langle k\rangle_{\ge2}$&2.8&3.4&4.1&4.9&5.6&6.5&7.4&8.4&9.4&10.4&11.5&12.7&13.8&14.8&16.0&17.1\\\bottomrule
\end{tabular}
\end{center}
Layer~1 is exact for $\relu$ of a centred Gaussian: $F/V=\frac14/(\frac12-\frac1{2\pi})=0.733$ and $\langle k\rangle_{\ge2}=2.75$.

\begin{proposition}[operator growth; measured]
As functions of the input, the neurons' first chaos decays to under a fifth of their variance. The mean degree of the rest grows linearly in
depth, about one degree per layer.
\end{proposition}
This is the classical counterpart of operator spreading in scrambling circuits, where the size of a Heisenberg operator grows linearly in
time. Its consequence for algorithm design is blunt. A degree-$d$ chaos component in $n=1024$ variables has about $n^d/d!$ coefficients. Any
expansion in a basis fixed in advance, whether Pauli-path or Heisenberg propagation in input Hermite coordinates, would need degree around
$17$ at the output. It is hopeless.

\subsection{Why closures work at all: frames that co-evolve with the weights}
Stage~15's wall-jet theorem expands each neuron's $\relu$ in the Appell (Wick) polynomials of \emph{its own layer's law}. In that frame the
births at one layer have low order: stage~14's kink law puts order $K$ at resolution radius $2\sqrt K$, and $K\le4$ suffices. Composing frames
telescopes the degree. Measured in the fixed input frame, the degree grows by one per layer; each frame change $\mu_l^{-1}\star\mathrm{id}$ absorbs that
growth, so in the moving frames it stays bounded.
\begin{itemize}[leftmargin=*]
\item \emph{The frame is a function of the instance.} The frame of layer $l$ is $(\mu_l,C_l)$, computed from $W_1,\dots,W_l$.
\item \emph{So the interdependence is forced.} The representation co-evolves with the sampled weights, and any tractable estimator must
have this property. The chain is one instance of it, not the template.
\item \emph{Quantum reading.} Classical simulations of random circuits by Pauli propagation truncate operator weight in a fixed basis. For
deep \textsc{ReLU} networks the truncation must be done in moving frames, which is the interaction picture with the Wick map as frame
change.
\end{itemize}
The open question this raises is about frame \emph{families}. Gaussian centred frames are one choice. Section~\ref{sec:gamma} shows
that reshaping a Gaussian frame variationally does not help; non-Gaussian frames are untested.

\section{Transport memory: holonomy, the normalizer, and lattice mismatch}
\label{sec:transport}
\subsection{The covariance chain does not determine the transport}
Fluctuations, and births, are carried from $z_{l-1}$ to $z_l$ by the first jet $B_l=W_l\diag\Phi(r_l)$. Whiten with the backbone covariances:
\[
\tilde O_l=C_l^{-1/2}B_lC_{l-1}^{1/2},\qquad \tilde O_l\tilde O_l^\top=I-C_l^{-1/2}\Sigma^{\rm birth}_lC_l^{-1/2},\qquad \Sigma^{\rm birth}_l=C_l-B_lC_{l-1}B_l^\top\succeq0 .
\]
So $\tilde O_l$ is a contraction times a rotation $U_l$. The contraction is the share of variance newly born from higher Mehler chaos. The
rotation is the network's frame turning, and $(C_{l-1},C_l)$ leaves it free: every $C_l^{1/2}RC_{l-1}^{-1/2}$ with $R\in O(n)$ is compatible with the
same pair of covariances. The holonomy note makes the same point sharply. Even the canonical symmetric-squeeze connection
$\dot T=\frac12\dot QQ^{-1}T$ depends on the \emph{path} of covariances, through the curvature $\frac14[HQ^{-1},KQ^{-1}]$, and a contact-preserving loop
returns with a rotation (checked, Section~\ref{sec:audit}).
\begin{proposition}[transport memory; proved]
Any estimator that propagates non-Gaussian content must carry the actual transports $B_l$ (or their action on its sources), and not
only the backbone $(\mu_l,C_l)$. Two instances with identical backbones but different rotations $U_l$ send the same births to different
places.
\end{proposition}

\subsection{The normalizer: why slices cannot be propagated as slices}
\begin{theorem}[hard-core normalizer, from the holonomy note]
An orthogonal $R$ preserves the square-free quadratic sector $\operatorname{span}\{z_iz_j:i<j\}$ if and only if it is a signed permutation.
\end{theorem}
\begin{corollary}[no closed diagonal closure; proved]
Suppose a closure keeps only neuron-diagonal data of the non-Gaussian state at each layer: the hub slices D3, D21 and $\kappa_4$
diagonals. It is closed under transport exactly when the transport, in the neuron bases, is a signed permutation. For Gaussian $W_l$ this
fails almost surely. Slices must therefore be \emph{regenerated} at every readout from sources carried with their legs, or from rows
transported backwards.
\end{corollary}
This is the exact form of stage~9's two empirical laws: Hadamard readouts defeat Frobenius truncation, and old slices correlate only
$0.3$--$0.5$ with young ones. A hub-localized leg $e_m$ becomes, after one transport, the dense column $\Phi_mW_{:,m}$, with inverse
participation ratio of order $3/n$. The chain's CP legs are one realization of transport memory. Stage~15's Heisenberg reads are the
other.

\subsection{Lattice reading}
$B_n=\mathrm{Aut}(\Z^n)\cap O(n)$.
\begin{itemize}[leftmargin=*]
\item Each layer's neuron basis defines a coordinate lattice $\Z^n$ in pre-activation space. Its reflections are the folds and its
orthants are the cells (stage~14: the layer map $z=Wx$ is Klartag's normalizing map).
\item The diagonal (hub) sector is the invariant sector of that lattice's automorphism group.
\item The transport between consecutive layers is a generic linear map, not an automorphism of the next layer's lattice.
\end{itemize}
So \emph{the memory obstruction is a lattice mismatch}: consecutive coordinate lattices are in general position. In the holonomy note's
example ($\Lambda=\Z e_1\oplus2\Z e_2$), contact-preserving deformations still produce rotations outside $\mathrm{Aut}(\Lambda)$, and irrational ones
produce a noncommutative torus. In the network the depth path is not a loop, but the same non-integrability makes the frame
unrecoverable from its endpoints.

\section{The carré du champ, Dirac energies, and two designs ruled out}
\label{sec:gamma}
\subsection{The product defect of the input heat flow is the network's covariance}
For the Ornstein--Uhlenbeck flow on inputs, $\Gamma(f,g)=\nabla f\cdot\nabla g$. The companion's accumulated product-defect identity, integrated to infinite
time, gives
\[
\Cov\big(f(X),g(X)\big)=\int_0^1\E\big[\nabla f(X)\cdot\nabla g(X_\rho)\big]\,d\rho,\qquad X_\rho=\rho X+\sqrt{1-\rho^2}\,X'.
\]
For neurons this is an integral over input overlaps of two-replica Jacobian correlators, $\Cov(h_l)=\int_0^1\E[J_l(X)J_l(X_\rho)^\top]\,d\rho$.
\emph{Check} on a small He network ($n=32$, $L=3$): the relative Frobenius difference from the Monte Carlo covariance is $0.008$, against
noise of about $0.004$, with diagonal ratio $0.998$.

The companion's Dirac energy, with coordinate derivations and $G=I$, is $\Tr C_\rho(A)=\E|\nabla A|^2$, the $G$ of Section~\ref{sec:growth}. Its excess over the
variance is the Poincaré deficit, $\E|\nabla h|^2-\Var h=\sum_k(k-1)\|h_k\|^2\ge0$: the higher-chaos (magic) energy of the neuron.

\subsection{Design ruled out: the leaf estimator with global Dirac energies}
For a bias-free \textsc{ReLU} network and Gaussian input, every output is $1$-homogeneous. Euler and Stein then give the exact leaf
decomposition
\[
\E f=\E\Delta f=\sum_{l,j}p_{z_{lj}}(0)\;\E\big[|\nabla z_{lj}|^2\,\partial f/\partial h_{lj}\;\big|\;z_{lj}=0\big],
\]
the transverse measure of each wall times the energy and sensitivity on it (stage~11). This offers a second estimator, independent of any
law closure. Decouple each wall term into
\begin{itemize}[leftmargin=*]
\item the backbone wall density $\varphi(r)/s$;
\item the \emph{global} Dirac energy $(W_lK_{l-1}W_l^\top)_{jj}$, from the Jacobian-Gram recursion
\[K_l=P_l\circ(W_lK_{l-1}W_l^\top),\]
with gate second moments $P_l$;
\item the first-jet cocycle as the sensitivity.
\end{itemize}
It costs about $2n^3$ per layer. On network~0 it is exact at layer~1 (MSE $7\times10^{-10}$, the truth's noise). At layer~2 its MSE is
$3.5\times10^{-2}$, against the Gaussian closure's $1.5\times10^{-7}$, and it grows to $8$ at layer~16.
\begin{itemize}[leftmargin=*]
\item \emph{Why it fails.} The global energy exceeds the variance of $z$ by the Poincaré deficit. The measured ratio of mean Dirac
energy to variance is $1.00$, $1.47$, $2.54$ and $13.7$ at layers $1$, $2$, $4$ and $16$, which is the chaos ladder of Section~\ref{sec:growth}.
\item \emph{Where the exact identity hides the cancellation.} In the exact identity this excess is cancelled between layers: a wall at
layer $l-1$ conditions the gates above it.
\item \emph{Consequence.} The leaf decomposition is exact but not term-wise stable, and the energy that enters the mean is the
\emph{transverse} one, $\E[|\nabla z|^2\mid z=0]$. In the companion's language, contact matrices must be evaluated in the leaf-conditioned state, not the
global state.
\end{itemize}
This also explains why stage~9's Euler--Stein merge had to use the chain's own Laplacian response rather than a wall sum.

\subsection{Design ruled out: minimal-sensitivity Wick frames}
The truncated Wick expansion of $\E\relu(Z-c)$ around a Gaussian frame $N(\kappa_1,v)$ depends on $v$; the exact answer does not. Principles of
minimal sensitivity pick a stationary $v$ instead of the centred $v=\kappa_2$. For $Z=g+a(g^2-1)$:
\begin{center}\small
\begin{tabular}{lcc}\toprule
case & centred frame error & stationary frames ($v/\kappa_2$: error)\\\midrule
$a=0.15$, wall $0.3$, order 3 & $+4.7\times10^{-3}$ & $1.22$: $+2.3\times10^{-3}$\\
$a=0.15$, wall $0.3$, order 4 & $-1.0\times10^{-2}$ & $0.69$: $-9.2\times10^{-3}$; $1.22$: $-1.1\times10^{-2}$\\
$a=0.15$, wall $-0.8$, order 4 & $+1.8\times10^{-3}$ & $0.51$: $-1.9\times10^{-2}$\\
$a=0.3$, wall $0.5$, order 4 & $-2.5\times10^{-2}$ & $1.41$: $-2.8\times10^{-2}$\\\bottomrule
\end{tabular}
\end{center}
Stationary frames are sometimes absent and are no better on balance. The centred frame, which is EFPTZ's unique centring character, is the
right choice of variance. Non-Gaussian frame \emph{families} remain untested.

\section{Interdependence of weights and estimate, made exact}
\label{sec:interdep}
The question was how an optimal system could fail to couple the sampled weights and its estimate, perhaps dynamically. There are four
exact couplings, ordered by how much they are forced.

\subsection{Forced: the frame}
Section~\ref{sec:growth} shows that the representation must be re-centred at every layer by the instance's own law. The frame of layer $l$ is a
function of $W_1,\dots,W_l$, and it changes as the computation proceeds.

\subsection{Free: the estimator learns its own correction from its own rows}
\begin{theorem}[row-exchangeable self-calibration; proved]\label{thm:selfcal}
Fix a layer $l$ and the upstream state $U$, meaning the weights below $l$ and the estimator's state up to layer $l-1$. For each row $w_i$ of $W_l$,
let $\pi_i=\pi(w_i;U)\in\R^p$ be a cheap per-row feature vector and $\beta_i=\beta(w_i;U)$ an expensive per-row quantity, for instance a higher-order birth.
Choose a uniformly random set $S$ of $q$ rows, fit $\hat a$ by least squares of $\beta$ on $\pi$ over $S$, and set $\hat\beta_i=\hat a\cdot\pi_i$ off $S$. Then:
\begin{enumerate}[leftmargin=*]
\item Given $U$, the pairs $(\pi_i,\beta_i)$ are i.i.d., because the rows are.
\item $\hat a$ converges to the population regression $a_*$, with excess error $O(p/q)$ relative to the residual.
\item The estimator's added mean squared error at layer $l$ is $(1-\frac qn)(1-R^2)\,\E\beta^2\,(1+O(p/q))$, where $R^2$ is the population fit.
\item By the error cocycle, this error reaches the output additively, multiplied by the layer's downstream gain.
\end{enumerate}
\end{theorem}
\begin{proof}
(1) The rows are i.i.d.\ and independent of $U$. (2) and (3) are least squares on i.i.d.\ samples. (4) is Theorem~3.2 of stage~15.
\end{proof}
In this precise sense an estimator learns its own correction from its own weights. The training set is the instance's rows, and what makes
it valid is exchangeability, not any independence of weights and estimate. Corollary~3.3 of stage~15 bounds $R^2$ by the isotropic share
whenever $\pi$ holds only per-neuron or field features. A large $R^2$ needs anisotropic proxies, such as a short-window or first-order
version of $\beta$ itself.

\emph{Where it should pay, and the test.} It pays where per-row costs dominate and such a proxy exists. Late-layer covariance births are
the case: layers $11$--$16$ hold $81\%$ of the exact chain's births, and $59\%$ of the last-layer birth is covariance.
\begin{itemize}[leftmargin=*]
\item The expensive quantity: the second-order covariance-birth readout of a row (the $(2,2)$, $(3,1)$ and $\kappa_3$-pair diagrams),
computed on $q\approx n/8$ rows.
\item The proxy: the first-order readout plus the field jets.
\item The prediction: the share $R^2$ of the channel is recovered at about an eighth of its readout cost.
\end{itemize}

\subsection{Free: geometry coupled to the readout}
The holonomy note's principle $\mathcal V(w)=\langle f,\mathsf K_w^{-1}f\rangle$ is convex in the weights $w$, and its optimality conditions equalize the edge
energies of the Poisson solution. Our analogue is stage~15's allocation:
\begin{itemize}[leftmargin=*]
\item The predicted MSE is convex in the capture fractions.
\item At the optimum the marginal MSE reduction per FLOP is equal across the (layer, row, diagram) choices used.
\item The Poisson solution is the adjoint: the downstream gains.
\end{itemize}
The two are coupled self-consistently. The allocation depends on the adjoint, and the adjoint depends on what was computed. Two passes,
or one pass with revision, solve it.

\subsection{Free: task-preserving instruments}
The companion's condition $\sum_\alpha\mathcal I_\alpha^*(A_\alpha)=A$ has a classical realization.
\begin{itemize}[leftmargin=*]
\item Branch on the sign of a chosen functional $u\cdot x$. By symmetry each branch has mass $\frac12$.
\item Run each branch with its own Wick frame. Then $\E f=\sum_\alpha P(\alpha)\E[f\mid\alpha]$ holds exactly.
\item A branch resolves all orders of non-Gaussianity along $u$, at twice the cost.
\end{itemize}
Its value depends on concentration, and the residual is not concentrated. In the experimental session, making the covariance exact in the top
$128$ collective modes gave only $16$--$32\%$, which a branching tree could approach only with $2^{128}$ leaves. Instruments are the right tool for
a coherent residual and the wrong one for an incoherent residual.

\subsection{The design space}
An estimator is a point in (frame family) $\times$ (partition blocks evaluated on the instance) $\times$ (rows) $\times$ (layers) $\times$ (instruments).
\begin{itemize}[leftmargin=*]
\item \emph{The standard chain.} Gaussian centred frames, pair and three-block diagrams, all rows, all layers, no instruments.
\item \emph{Forced by this stage.} Co-evolving frames and transport memory.
\item \emph{Ruled out here.} Leaf estimators with global energies, and minimal-sensitivity Gaussian frames.
\item \emph{Open, with protocols.} Non-Gaussian frame families, diagram selection by annealed two-copy variances,
self-calibration over rows, and late-layer allocation.
\end{itemize}

\section{Cross-validation with the experimental session}
\label{sec:crossval}
The experimental session tested stage~13 on v56 and on the truth of network~0, in §9 of its
\texttt{notes/\allowbreak fresh-weight}. Its results read cleanly in
this stage's terms.
\begin{center}\small
\begin{tabular}{>{\raggedright\arraybackslash}p{6.2cm}>{\raggedright\arraybackslash}p{7.8cm}}\toprule
experimental finding & reading\\\midrule
gain saturates; increments uncorrelated ($V_l/4\approx0.0082$ flat from layer~13) & stage~12--13 prediction (saturation, not mean reversion)\\
the chain's slices carry the gain shape (82--94\% for $\kappa_3$, K22, $\kappa_4$ diagonal) & the Brauer (dilation) component of each slice (Proposition~\ref{prop:annealed})\\
setting slice gain amplitudes to the truth's: $+56\%$, $+50\%$ (predicted $-10$ to $-30\%$) & the counterterms already hold the coherent sector; coherent errors add in amplitude across layers (stage~15), so local amplitude changes are amplified\\
final error has $0.4$--$3.7\%$ of its energy along the mean & stage~15: collective share $0$--$3\%$ for the exact chain; the residual is incoherent births\\
gain package with the old-tier ranks cut recovers only $5.3\%$ & the lost old-tier content is three-block (stars), purely quenched; a Brauer package cannot replace it\\
top-32 / top-128 covariance modes exact: $-12$ to $-16\%$ / $-16$ to $-32\%$; eigenvalues only: $+9\%$ & row theorem: a row's covariance birth $\frac12\J_2\,w^\top\delta C\,w$ weighs every mode by its share of $\|\delta C\|_F^2$; eigenvalues alone leave the anisotropic error\\
layers 0--6 create a few percent of the final error & stage~15 measured $5\%$ (layers $1$--$6$)\\\bottomrule
\end{tabular}
\end{center}
\paragraph{What this suggests for the experimental session.}
\begin{enumerate}[leftmargin=*]
\item \emph{Leave the coherent sector alone.} It is held by the counterterms. Further work on it is fighting the amplification
mechanism.
\item \emph{Go after the late-layer covariance births, full rank.} These are the $59\%$ channel. Use the consistent second-order
diagram set, the $(2,2)$, $(3,1)$ and $\kappa_3$-pair terms of the wall-jet expansion at matching order, in layers $11$--$16$. Fund it from layers
$1$--$6$, and refit the counterterms afterwards.
\item \emph{Decide self-calibration with one cheap diagnostic first.} At layers $12$--$16$ on network~0, measure the $R^2$ of the per-row
second-order covariance birth on its first-order proxy and the field jets. If $R^2\gtrsim0.7$, Theorem~\ref{thm:selfcal} cuts the readout cost of the
channel about eightfold.
\end{enumerate}

\section{Checks}
\label{sec:verify}
\texttt{scripts/\allowbreak verify\_\allowbreak stage16.py} runs in about six minutes. Most of that is the Jacobian scan of network~0; set
\texttt{SKIP\_SLOW=1} to skip it. The transcript is \texttt{notes/\allowbreak stage16/\allowbreak verify\_\allowbreak stage16.txt}.
\begin{center}\small
\begin{tabular}{>{\raggedright\arraybackslash}p{5.4cm}>{\raggedright\arraybackslash}p{8.6cm}}\toprule
statement & result\\\midrule
spin branching $p_\pm$ & $(j+1)/(2j+1)$, $j/(2j+1)$, sum $1$ at $j=\frac12,1,\frac72$\\
three-spin logical qubit & $S_{12}=-Z$, $S_{23}=\frac12Z+\frac{\sqrt3}2X$, $[S_{12},S_{23}]=-i\sqrt3Y$ exactly\\
Zhou v1 multi-time display & $\E[Z_a^2Z_b^2]$: $3.021$ (MC) vs $ab+2a^2+a=3.010$; product of means $1.33$\\
isotropic angular reference $G_1$ & $0.5$ and $0.0998$--$0.1004$ (pred.\ $0.1$), off-diagonal $\le3\times10^{-4}$\\
squeeze holonomy of a contact loop & angle$/(a^2/4)=0.985,0.965,0.935$ at $a=0.02,0.05,0.1$; orthogonal to $5\times10^{-13}$\\
hard-core normalizer & $z_i^2$ leak $0$ (signed permutation), $0.39$ (random rotation)\\
mixture $\kappa_4=\Cov(Q)$ & $1.573/1.558$, $0.538/0.518$, $-0.007/-0.001$\\
Bernoulli realization relation & $:\!b^2\!:-:\!b\!:\;=-2p(B-p)$ exactly\\
commutant tower & $15,4,3.01$ ($2k=4$) and $203,31,15.17$ ($2k=6$) for $S_n,B_n,O(n)$\\
replica covariance identity & relative Frobenius $0.008$ (noise $\approx0.004$), diagonal ratio $0.998$\\
chaos ladder (network 0) & $F/V$: $0.73\to0.19$; $\langle k\rangle_{\ge2}$: $2.76\to17.14$ over layers $1\to16$\\
isotropic share of truth $-$ Gaussian & $0.60$--$0.66$; exact chain captures $0.991$--$0.997$\\
minimal-sensitivity frames & no gain over the centred frame (table in Section~\ref{sec:gamma})\\
wall-sum estimator & MSE $3.5\times10^{-2}$ (layer 2), $8.3$ (layer 16); Dirac/variance $1.47$, $13.7$\\\bottomrule
\end{tabular}
\end{center}
Not checked numerically:
\begin{itemize}[leftmargin=*]
\item the self-calibration theorem, whose protocol is in Section~\ref{sec:crossval};
\item the claim that three-block content is purely quenched, which follows from Proposition~\ref{prop:annealed}(1), beyond the vanishing of annealed
odd cumulants of pre-activations.
\end{itemize}

\end{document}
